% =====================================================================
% Section 3 — Methodology
% =====================================================================

\section{Methodology}\label{sec:methodology}

\subsection{System under study}\label{sec:method-system}

The system under study is the valve train of a high-performance cylinder
head conceived for the Volkswagen AP 1.8 platform. The base engine has a
bore of 81.0\,mm, a stroke of 86.4\,mm, a total displacement of
1\,781\,cm$^{3}$, and a four-cylinder in-line architecture
\citep[][\S~1.1]{porto2026dissertation}. The cylinder head is of the
pent-roof type, with inclined valves, asymmetric quench, gasoline direct
injection (GDI), and an iridium spark plug; these features establish the
mechanical and thermal context of the valve train but are not the
subject of the risk analysis presented here. The reader is referred to
\citet{porto2026dissertation} for the full description of the cylinder
head and the air-flow optimisation.

\subsection{Mechanical and operational premises}
\label{sec:method-premises}

The premises adopted for the present risk analysis are summarised in
Table~\ref{tbl:premises} and follow the design choices of the
dissertation~\citep{porto2026dissertation} and the associated sizing
notes for the intake valve and the valve springs.

\begin{table*}[!htbp]
\caption{Mechanical and operational premises adopted for the risk analysis of the AP 1.8 valve train. Values taken from the dissertation~\citep{porto2026dissertation} and the associated sizing notes for the intake valve and the valve springs.}\label{tbl:premises}
\centering
\resizebox{\textwidth}{!}{%
\begin{tabular}{@{}lllll@{}}
\toprule
\multicolumn{2}{l}{\emph{Engine and operating regime}} & & \multicolumn{2}{l}{\emph{Dual valve-spring assembly (Cr--Si steel, ASTM A877/A877M)}} \\
\midrule
Bore & 81.0~mm & & Outer spring stiffness $k_{\mathrm{out}}$ & 38.06~N\,mm$^{-1}$ \\
Stroke & 86.4~mm & & Inner spring stiffness $k_{\mathrm{in}}$ & 61.43~N\,mm$^{-1}$ \\
Total displacement & 1\,781~cm$^{3}$ & & Combined stiffness $k_{\mathrm{total}}$ & 99.49~N\,mm$^{-1}$ \\
Cylinders & 4 in line & & Seat preload $F_{\mathrm{seat}}$ & 780~N \\
Maximum engine speed & 10\,000~rpm & & Force at full lift $F_{\mathrm{open}}$ & 1\,985.8~N \\
Camshaft fund.\ freq.\ $f_{\mathrm{cam,max}}$ & 83.3~Hz & & Outer spring solid height & 33.6~mm \\
Compression ratio & 14.5\,:\,1 & & Inner spring solid height & 32.0~mm \\
\addlinespace
\multicolumn{2}{l}{\emph{Valves, rocker and cam}} & & Outer spring height at full lift & 34.38~mm \\
\cmidrule(r){1-2}
Intake valve material & Ti--6Al--4V & & Total spring assembly mass & 152.72~g \\
Exhaust valve material & INCONEL 751 & & Effective spring mass ($m_{\mathrm{molas}}/3$) & 50.91~g \\
Roller rocker / retainer / keeper & 7075--T6 aluminium & & & \\
Camshaft material & SAE 8620 case-hardened & & & \\
\cmidrule(r){4-5}
Maximum valve lift & 12.12~mm & & \multicolumn{2}{l}{\emph{Direct translational mass (intake side)}} \\
Cam accel.\ peak (positive) & 14\,544~m\,s$^{-2}$ & & Valve + retainer + keeper & 43.48~g \\
Cam decel.\ peak (critical) & 20\,362~m\,s$^{-2}$ & & With reflected rocker & 49.59~g \\
\cmidrule(r){4-5}
& & & \multicolumn{2}{l}{\emph{Direct translational mass (exhaust side)}} \\
& & & Valve + retainer + keeper & 54.25~g \\
& & & With reflected rocker & 60.36~g \\
\bottomrule
\end{tabular}}
\end{table*}

\subsection{Construction of the qualitative BowTie}
\label{sec:method-bowtie-qual}

The BowTie diagram is constructed around the top event \emph{loss of
follower contact between valve and cam}. This formulation
intentionally aggregates the three failure modes of
Section~\ref{sec:bg-failure-modes} that share a single mechanical
signature: valve float, valve bounce, and approach to coil bind. The
remaining modes (resonance and fatigue) act as enabling threats rather
than as the top event itself.

The threats placed on the left side of the bow are derived from the
classical literature \citep{heywood1988,taylor1985,heisler2002} and
from the sizing constraints of the AP 1.8 head. They are: (i)
insufficient spring stiffness for the inertial regime; (ii) excessive
moving mass; (iii) aggressive cam profile; (iv) spring heating in
operation; (v) out-of-tolerance thermal clearance; and (vi) resonance
between cam harmonics and spring natural frequency. The preventive
barriers placed between threats and the top event are: (i) dynamic
sizing of the spring relative to the cam harmonics; (ii) dual or
nested springs; (iii) low-mass tappets; (iv) high-fatigue-resistance
spring material; and (v) modal analysis and dynamometer validation.

The consequences placed on the right side of the bow are: (i)
valve--piston collision; (ii) catastrophic cylinder-head failure; (iii)
immediate loss of compression; (iv) oil contamination by debris; and (v)
unscheduled engine shutdown. The mitigative barriers between the top
event and these consequences are: (i) electronic rev limiter; (ii)
cam-position sensor; (iii) operating envelope defined at design stage;
(iv) torque or fuel cut-off; and (v) predictive maintenance via
vibration signals.

\subsection{Quantitative coupling layer}
\label{sec:method-bowtie-quant}

To overcome the qualitative limitation of the classical BowTie, each of
the five preventive barriers is associated with a deterministic
indicator and a pass criterion, summarised in
Table~\ref{tbl:quant-criteria}. The indicators are computed from the
analytical sizing of the intake valve and of the valve springs,
which is reported in the underlying dissertation
\citep{porto2026dissertation} and in two associated sizing notes
on the intake valve and the spring.

\begin{table*}[!htbp]
\caption{Quantitative indicators associated with the preventive barriers of the BowTie diagram. The pass criteria follow conventional design rules in the internal-combustion-engine literature and the project specification of the AP 1.8 sizing.}\label{tbl:quant-criteria}
\centering
\resizebox{\textwidth}{!}{%
\begin{tabular}{@{}lll@{}}
\toprule
Barrier & Indicator & Pass criterion \\
\midrule
Dynamic follower contact & $R = F_{\mathrm{spring,crit}}/F_{\mathrm{inertia,crit}}$ & $R \geq 1.20$ (project) \\
Spring natural frequency & $f_{n}=\tfrac{1}{2}\sqrt{k/m_{s}}$ & $f_{n}/f_{\mathrm{cam,max}}\geq 4$~\citep{heywood1988} \\
Coil-bind margin (outer) & $M = H_{\mathrm{open}}-H_{\mathrm{cb}}$ & $M \geq 1.5$~mm at 10\,000~rpm \\
Fatigue (Goodman) & $1/n_{f}=\tau_{a}/S_{se}+\tau_{m}/S_{su}$ & $n_{f}\geq 1$~\citep{shigley2020} \\
Hot clearance$^{\dagger}$ & $c(T_{\max})$ & $c(T_{\max})\geq 0$ over the full map \\
\bottomrule
\end{tabular}}\\
{\footnotesize $^{\dagger}$Hot clearance is treated in this paper as a \emph{qualitatively verified} barrier outside the deterministic layer; only a first-order order-of-magnitude estimate is provided in Section~\ref{sec:res-quant}. A complete differential thermal model is identified as a priority direction for future work in Section~\ref{sec:disc-limitations}.}
\end{table*}

The dynamic-ratio criterion $R\geq 1.20$ corresponds to the project
specification adopted in the AP 1.8 sizing and is consistent with
the design practice for high-speed valve trains
\citep{heywood1988,taylor1985,heisler2002}.

The natural-frequency criterion follows the rule of thumb that the
first surge mode of the spring should sit above the dominant cam
harmonics \citep{heywood1988}. Two equivalent formulations are
available in the literature: the one-degree-of-freedom approximation
$f_{n} = \tfrac{1}{2}\sqrt{k/m_{s}}$ for a coil spring with both ends
fixed \citep[\S~10-26]{shigley2020}, and the continuous-medium
expression
$f_{n} = (d/(2\pi N D_{m}^{2}))\sqrt{G/(2\rho)}$
\citep[\S~6.6]{heywood1988}, which reduce algebraically to the same
result when $m_{s}$ in the first form is taken as the active-coil
mass. In the present analysis the conservative form
$f_{n} = \tfrac{1}{2}\sqrt{k/m_{s,\mathrm{total}}}$ is adopted with the
\emph{total} spring mass (including end coils), which yields the
lower bound of the surge frequency and is the more demanding test for
the rule-of-thumb threshold.

The coil-bind margin of 1.5--2.0~mm is the
practical specification adopted in the AP 1.8 sizing to absorb
manufacturing tolerances, mounting variations, and dynamic excitations
typical of operation at 10\,000~rpm. The fatigue criterion uses the Goodman alternating-mean factor of
safety (Eq.~\ref{eq:goodman}) with the shear endurance limit and the
ultimate shear strength of ASTM A877/A877M chrome--silicon
valve-spring steel after shot peening and preset
\citep{shigley2020,xing2021valvespringEAC,ko2022springfatigueflaws}.
The criteria themselves are drawn from established design practice; the
contribution of the QC-BowTie is not a new criterion but the disciplined
coupling of each barrier to a measurable indicator and the propagation of
that indicator to a compliance-failure probability
(Section~\ref{sec:res-mc}).

\paragraph{Status of the pass criteria.} Two of the five criteria are
external standards, the surge-frequency margin
$f_{n}/f_{\mathrm{cam,max}}\geq 4$ \citep{heywood1988} and the Goodman
fatigue criterion $n_{f}\geq 1$ \citep{shigley2020}. The dynamic-ratio
target $R\geq 1.20$ and the coil-bind target $M\geq 1.5$~mm are
\emph{project specifications} adopted in the AP 1.8 sizing rather than
universal thresholds. A barrier reported below one of these targets is
therefore a margin deficit \emph{relative to the adopted target}, not an
absolute prediction of failure; to keep this distinction auditable, every
barrier in Table~\ref{tbl:quant-results} is reported together with its
underlying quantitative margin, so that a reader adopting a different
target can re-read the classification directly. The influence of the two
speed-sensitive targets is examined explicitly by the operating-envelope
sweep of Section~\ref{sec:res-envelope}.

\subsection{Cross-check procedure}\label{sec:method-cross-check}

For the AP 1.8 head, the values of all indicators in
Table~\ref{tbl:quant-criteria} are computed under the operating envelope
of Table~\ref{tbl:premises}. A barrier is declared \emph{active} if the
indicator satisfies the pass criterion; otherwise it is declared
\emph{at risk}. The single barrier with the smallest absolute margin is
declared the \emph{critical barrier}: it is the first that would fail
under a parametric drift of the inputs (mass increase, stiffness loss
through heating, harder cam profile, etc.). A tornado plot of the
sensitivity of each indicator to its dominant input is reported in
Section~\ref{sec:results} (Fig.~\ref{fig:tornado}) to make the
identification of the critical barrier transparent.

The qualitative BowTie of
Section~\ref{sec:method-bowtie-qual} and the quantitative layer of
Section~\ref{sec:method-bowtie-quant} are then cross-checked against a
conventional FMEA/RPN ranking of the same threats. The comparison is
discussed in Section~\ref{sec:disc-fmea-comparison}. The complete
QC-BowTie pipeline is summarised in Fig.~\ref{fig:flow}.

\begin{figure}[!htbp]
 \centering
 \resizebox{\linewidth}{!}{% Method-flow figure for the QC-BowTie pipeline.
\begin{tikzpicture}[
    node distance=7mm and 8mm,
    box/.style={draw, rounded corners, align=center, font=\footnotesize,
                inner sep=3.5pt, text width=26mm, minimum height=12mm},
    arr/.style={-{Latex[length=2.2mm]}, thick},
]
\node[box, fill=blue!6] (b1) {Qualitative BowTie\\ (top event, threats,\\ preventive barriers)};
\node[box, fill=blue!6, right=of b1] (b2) {Couple each barrier to a\\ deterministic indicator\\ $+$ pass criterion};
\node[box, fill=green!8, right=of b2] (b3) {Deterministic\\ evaluation:\\ active / at-risk};
\node[box, fill=orange!10, below=of b3] (b4) {Monte Carlo\\ (ASTM A877\\ tolerances)};
\node[box, fill=orange!10, left=of b4] (b5) {Barrier status as\\ $P[\mathrm{fail}]$: active /\\ marginal / at-risk};
\node[box, fill=gray!10, left=of b5] (b6) {Contrast with\\ FMEA / RPN};
\draw[arr] (b1) -- (b2);
\draw[arr] (b2) -- (b3);
\draw[arr] (b3) -- (b4);
\draw[arr] (b4) -- (b5);
\draw[arr] (b5) -- (b6);
\node[box, fill=red!6, above=of b2, text width=40mm] (cv)
    {Construct-validity check vs.\ four published failure cases};
\draw[arr, dashed] (cv) -- (b2);
\end{tikzpicture}
}
 \caption{The QC-BowTie pipeline: the qualitative BowTie is enriched by coupling each preventive barrier to a deterministic indicator with a pass criterion; the indicators are evaluated deterministically and then propagated by Monte Carlo under manufacturing tolerances to a per-barrier compliance-failure probability $P[\mathrm{fail}]$ (active/marginal/at-risk), and the outcome is contrasted with a conventional FMEA/RPN ranking. The preventive-barrier set is informed by a construct-validity check against four published failure cases.}\label{fig:flow}
\end{figure}
